%%%%%%%%%%%%%%%%%%%%%%%%%%%%%%%%%%%%%%%%%%%%%%%%%%%%%%%%%%%%%%%%%
% 							SETTINGS 							%
%%%%%%%%%%%%%%%%%%%%%%%%%%%%%%%%%%%%%%%%%%%%%%%%%%%%%%%%%%%%%%%%%

\documentclass[12pt]{article}

\usepackage{
	amsmath,
	amssymb,
	geometry,
	graphicx,
	esint
}


\newcommand{\proptoinverse}{\mathrel{\mskip1mu\reflectbox{$\propto$}\mskip-1mu}}

\geometry{margin=1in}

\author{Álvaro Fernández Barrero and Vicente Gabriel Piñar Ojeda}
\title{Electromagnetism}
\date{04-10-2026}

\begin{document}
	\maketitle
	
	%%%%%%%%%%%%%%%%%%%%%%%%%%%%%%%%%%%%%%%%%%%%%%%%%%%%%%%%%%%%%%%%%%%%%
	% 							INTRODUCTION 							%
	%%%%%%%%%%%%%%%%%%%%%%%%%%%%%%%%%%%%%%%%%%%%%%%%%%%%%%%%%%%%%%%%%%%%%
	
	% Author: Álvaro Fernández Barrero
	% Date: 2026-10-5
	\section{Particle's charge}
	
	Every particle has an intrinsic property called charge. This is a numerical value $Q$ which can be positive or negative. Some classic examples are electrons with negative charge and protons with positive charge, which is the reason why, chemistry-wise, we use $e^-$ to refer to electrons and $p^+$ to refer to protons.
	
	The charge has a meaningful behavior: opposite charges (positive and negative) attract each other and same charges (positive and positive, or negative and negative) repel each other.
	
	This is key to understand everything that comes next as it is the fundamental idea of the whole topic.
	
	These charges have the unit coulomb ($[Q] = C$).
	
	%%%%%%%%%%%%%%%%%%%%%%%%%%%%%%%%%%%%%%%%%%%%%%%%%%%%%%%%%%%%%%%%%%%%%
	% 							COULOMB'S LAW 							%
	%%%%%%%%%%%%%%%%%%%%%%%%%%%%%%%%%%%%%%%%%%%%%%%%%%%%%%%%%%%%%%%%%%%%%
	
	% Author: Álvaro Fernández Barrero
	% Date: 2026-10-5
	\section{Coulomb's law}
	
	When we take a look to when two particles are close to each other and they are either repelling or attracting each other, we can try to find out a formula that describe the force acting upon the particles. By doing some experimentation, we see that:
	
	\[
		F \propto Q_1 \quad\quad F \propto Q_2 \quad\quad F \proptoinverse r^2
	\]
	
	Where $r$ is the distance between the two particles and $F$ the magnitude of the vector force $F = \|\vec{F}\|$.
	
	At this point, these proportions are pretty much the same as the gravitational force and, therefore, we will obtain a similar formula.
	
	Mixing these pieces together, we have:
	
	\[
		F = k \frac{Q_1 Q_2}{r^2}
	\]
	
	And again, by doing some experiments, we can obtain the value of the constant $k$:
	
	\[
		k = 8.99 \cdot 10^9 \frac{N \cdot m^2}{C^2} = \frac{1}{4 \pi \varepsilon_0} \frac{N \cdot m^2}{C^2}
	\]
	\[
		\therefore F = \frac{Q_1 Q_2}{4 \pi \varepsilon_0 r^2}
		\implies \vec{F} = \frac{Q_1 Q_2}{4 \pi \varepsilon_0 r^2} \hat{r}
	\]
	
	Where $\hat{r}$ is a unit vector pointing from one charge to the other.
	
	%%%%%%%%%%%%%%%%%%%%%%%%%%%%%%%%%%%%%%%%%%%%%%%%%%%%%%%%%%%%%%%%%%%%%
	% 							ELECTRIC FIELD 							%
	%%%%%%%%%%%%%%%%%%%%%%%%%%%%%%%%%%%%%%%%%%%%%%%%%%%%%%%%%%%%%%%%%%%%%
	
	% Author: Álvaro Fernández Barrero
	% Date: 2026-10-5
	\section{Electric field}
	
	Coulomb's law only work for two given charges, but the universe is full of them. That means that somehow we could attempt to have a more general formula for $n$ particles in our system. To do that we can think that a given particle with charge $Q$ is being constantly either pushed or pulled by different particles with different charges and, as this motion is provoked by individual forces, we can sum all of them up to have the total force applied to the particle with $Q$ charge.
	
	\[
		\vec{F} = \sum_{\mu = 1}^n \frac{Q_\mu Q}{4 \pi \varepsilon_0} \hat{r}_\mu
	\]
	
	Now that we have a nice expression for the forced applied to a particle given the $n$ charges in the system, we can perform some mathematical tricks to make it a little better to compute.
	
	\[
		\vec{F}
		= \frac{Q}{4 \pi \varepsilon_0} \sum_{\mu = 1}^n \frac{Q_\mu}{r_\mu^2} \hat{r}_\mu
	\]
	
	As $r = \|\vec{r}_\mu - \vec{r}\|$ by definition, being $\vec{r}$ the testing particle's position and $\vec{r}_\mu$ the $\mu$-charged particle's position, then:
	
	\[
		\vec{F}
		= \frac{Q}{4 \pi \varepsilon_0} \sum_{\mu = 1}^n \frac{Q_\mu}{\|\vec{r}_\mu - \vec{r}\|^2} \hat{r}_\mu
	\]
	
	And again, by how we have defined the parameters:
	
	\[
		\hat{r}_\mu = \frac{\vec{r}_\mu - \vec{r}}{\|\vec{r}_\mu - \vec{r}\|}
	\]
	
	Thus:
	
	\[
		\vec{F}
		= \frac{Q}{4 \pi \varepsilon_0} \sum_{\mu = 1}^n \frac{Q_\mu}{\|\vec{r}_\mu - \vec{r}\|^3} (\vec{r}_\mu - \vec{r})
	\]
	
	That is the formula after all this work, but the purpose of it all was actually to have a formula such that for any particle with any charge we can compute the force that is going to undergo. For the moment, that testing particle's charge is part of the formula, even though that is not something we want. Therefore, let us define a vector field $\vec{E}(\vec{r})$ as:
	
	\[
		\vec{E}(\vec{r}) = \frac{\vec{F}}{Q} = \frac{1}{4 \pi \varepsilon_0} \sum_{\mu = 1}^n \frac{Q_\mu}{\|\vec{r}_\mu - \vec{r}\|^3} (\vec{r}_\mu - \vec{r})
	\]
	
	This is what we call \emph{Electric field} and returns the force a particle would have given a position in the system according to the set particles' charge.
	
	By that, something easy to state is that, a particle with charge $Q$ will have a force at $\vec{r}$ of $Q\vec{E}(\vec{r})$ in the system.
	
	% Author: Álvaro Fernández Barrero
	% Date: 2026-10-9
	\subsection{Gauß's law}
	
	From here, the mathematics used are going to increase significantly using vector calculus, but the results are gorgeous and riveting.
	
	Let us get some more realistic and be in a 3D space, like in our daily life, which will allow us to use the Gauß's divergence theore.
	
	Something interesting happens when we want to compute the flux of the electric field over a surface $\partial V$, let us go through it.
	
	\[
		\Phi_E = \oiint_{\partial V} \vec{E} \cdot d\vec{S}
		= \oiint_{\partial V} \frac{1}{4 \pi \varepsilon_0} \sum_{\mu = 0}^{n} \frac{Q_\mu}{\|\vec{r}-\vec{r_\mu}\|^3}(\vec{r} - \vec{r_\mu}) \cdot d\vec{S}
	\]
	
	It is important to remark here that every charge $Q_\mu$ is \textbf{within} the surface $\partial V$.
	Then, we can take the constants out of the integral and use the linearity of the integrals to obtain:
	
	\[
		\Phi_E
		= \frac{1}{4 \pi \varepsilon_0} \sum_{\mu = 0}^{n} \oiint_{\partial V} \frac{Q_\mu}{\|\vec{r}-\vec{r_\mu}\|^3}(\vec{r} - \vec{r_\mu}) \cdot d\vec{S}
	\]
	
	To make it simple, let us assume $\partial V$ be a sphere.
	Now, we can place $Q_\mu$ outside the integral since it is a constant for it and express this integral in spherical coordinates. To make it simple, we will say that $\vec{r} - \vec{r_\mu} = \vec{r}_{p\mu}$ and $\|\vec{r} - \vec{r_\mu}\| = \rho$.
	
	Here, the differential $d\vec{S} = dS\hat{r} = \rho^2 \sin(\theta) d\theta d\varphi \hat{r}$ where $\rho$ is the radius of the sphere. Thus:
	
	\[
		\begin{aligned}
			& \Phi_E
			= \frac{1}{4 \pi \varepsilon_0} \sum_{\mu = 0}^{n} Q_\mu \oiint_{\partial V} \frac{1}{\rho^2} \hat{r}_{p\mu} \cdot (\rho^2 \sin(\theta) d\theta d\varphi \hat{r})
			= \frac{1}{4 \pi \varepsilon_0} \sum_{\mu = 0}^{n} Q_\mu \oiint_{\partial V} \frac{1}{\rho^2} \rho^2 \sin(\theta) d\theta d\varphi (\hat{r}_{p\mu} \cdot \hat{r})
			\\
			& \hspace{1.5em} = \frac{1}{4 \pi \varepsilon_0} \sum_{\mu = 0}^{n} Q_\mu \oiint_{\partial V} \sin(\theta) d\theta d\varphi (\hat{r}_{p\mu} \cdot \hat{r})
		\end{aligned}
	\]
	
	Here, $\hat{r}$ and $\hat{r}_{p\mu}$ must be two unit vectors pointing in the same direction, which means that $\hat{r} \cdot \hat{r}_{p\mu} = 1$:
	
	\[
		= \frac{1}{4 \pi \varepsilon_0} \sum_{\mu = 0}^{n} Q_\mu \oiint_{\partial V} \sin(\theta) d\theta d\varphi
	\]
	
	We know rewrite the surface $\partial V$ in terms of spherical coordinates which is $\partial V = [0, 2\pi] \times [0, \pi]$:
	
	\[
		\begin{aligned}
			& = \frac{1}{4 \pi \varepsilon_0} \sum_{\mu = 0}^{n} Q_\mu \int_{0}^{2\pi} \int_{0}^{\pi} \sin(\theta) d\theta d\varphi
			= \frac{1}{4 \pi \varepsilon_0} \sum_{\mu = 0}^{n} Q_\mu \int_{0}^{2\pi} [-\cos(\theta)]_{\theta = 0}^{\theta = \pi} d\varphi
			= \frac{1}{4 \pi \varepsilon_0} \sum_{\mu = 0}^{n} Q_\mu \int_{0}^{2\pi} 2 d\varphi
			\\
			& = \frac{1}{4 \pi \varepsilon_0} \sum_{\mu = 0}^{n} 2Q_\mu \varphi|_{\varphi = 0}^{\varphi = 2\pi}
			= \frac{1}{4 \pi \varepsilon_0} \sum_{\mu = 0}^{n} Q_\mu 4\pi
			= \frac{1}{4 \pi \varepsilon_0} 4\pi  
			= \frac{1}{\varepsilon_0} \sum_{\mu = 0}^{n} Q_\mu
		\end{aligned}
	\]
	
	Hence:
	
	\[
		\Phi_E
		= \oiint_{\partial V} \vec{E} \cdot d\vec{S}
		= \frac{Q_{in}}{\varepsilon_0},
		\hspace{2em}
		Q_{in} = \sum_{\mu = 0}^n Q_\mu
	\]
	
	In the result, we do not see that the radius of the sphere we have been using to integrate over the field is present. This is super interesting because it means that the result does not rely on it whatsoever: it does not matter the radius of the sphere we use, the result will be always the same.
	
	This is the first Maxwell's law, 4 laws that describe how electric and magnetic fields interact with each other, but this is a topic we will discuss later.
	
	With this result, we can obtain the differential form of it by using \emph{Gauß's divergence theorem} and thinking that $Q_{in}$ is actually that charge density $\rho$ within the given surface $\partial V$.
	
	\[
		Q_{in} = \iiint_{V} \rho dV
		\oiint_{\partial V} \vec{E} \cdot d\vec{S}
		= \iiint_V \nabla \cdot \vec{E} dV
	\]
	
	Thus, we can now use these forms in our original flux $\Phi_E$:
	
	\[
		\iiint_V \nabla \cdot \vec{E} dV
		= \iiint_V \frac{\rho}{\varepsilon_0} dV
	\]
	
	Here, since we are integrating over the same volume, the two integrals will have the same result when the integrating parts are the same, thus:
	
	\[
		\therefore \nabla \cdot \vec{E} = \frac{\rho}{\varepsilon_0}
	\]
	
\end{document}